% !TEX root = ../main.tex
% !TEX program = lualatex
\documentclass[../main.tex]{subfiles}
\IfSubfilesClassLoaded{\externaldocument{\subfix{../build/main}}}{}
% ====================
% File: 17K_jcids.tex
% ====================
\begin{document}
\IfSubfilesClassLoaded{\chapter{Joint Force Requirements Process}}{}
% === JFRP (EDO 3.3.3); filename retained for historical source links ===
\section{Joint Force Requirements Process}
\label{sec:jfrp}
\note{Current as of 2026-10-20. Coursebook 3.3.3, Version 7.0 (30 September 2026), teaches the JFRP as the current process. Its slides 33--35 reproduce excerpts labeled CJCSM 5123.01A, 5 August 2026; those excerpts are the source basis for the resourcing discussion below. Legacy JCIDS terminology is retained only for historical crosswalks~\autocite{edo-3-3-3-jfrp-2026}.}

\subsection{Summary}
\begin{description}
  \item[\textbf{Decision support triad.}] \ac{jfrp} defines requirements, \ac{ppbe} provides resources, and the \ac{was}, through the \ac{aaf}, delivers materiel capabilities. These systems operate continuously and must stay aligned~\autocite{edo-3-3-3-jfrp-2026}.
  \item[\textbf{Replacement.}] The Joint Staff rescinded \ac{jcids} on 15 January 2026. The Joint Staff catalogs corroborate CJCSI~5123.01J (15~January~2026, Change~1 on 5~August) and CJCSM~5123.01A (5~August~2026). Detailed process explanations below remain grounded in the supplied coursebook excerpts~\autocite{JCS-Requirements-Instructions-20261020,JCS-Requirements-Manuals-20261020,edo-3-3-3-jfrp-2026}.
  \item[\textbf{Service authority.}] Services and Components approve and certify their own requirements, sharing approved documents and program data for joint awareness and integration rather than waiting for routine joint validation~\autocite[slides 7--10, 34]{edo-3-3-3-jfrp-2026}.
  \item[\textbf{Joint focus.}] The \ac{jroc} concentrates on Joint Force Design, Joint Capability Integration, and \ac{ccmd} requirements; it annually prioritizes \acp{jop}~\autocite[slides 9, 30--33]{edo-3-3-3-jfrp-2026}.
  \item[\textbf{Documents.}] An initial \ac{srd} or \ac{icd} describes the gap before \ac{mdd}; an updated \ac{srd} or \ac{cdd} defines measurable capabilities before Milestone~B. Services may tailor document names and approval processes~\autocite[slides 7, 20--23]{edo-3-3-3-jfrp-2026}.
  \item[\textbf{Urgent needs.}] The \ac{jicr} consolidates the former \ac{juon}/\ac{jeon} pathways; distinguish urgent requirements approval from the acquisition pathway that delivers the solution~\autocite{edo-3-3-3-jfrp-2026,edo-3-3-2-acq-milestones-phases-2025}.
\end{description}

\subsection{Why JFRP Replaced JCIDS}
The reform seeks faster capability delivery by reducing sequential staffing, empowering Services, preserving the independent warfighter voice, and connecting strategy, analysis, requirements, and budgets. Component autonomy does not remove joint interoperability, analytic due diligence, or acquisition decision requirements~\autocite[slides 7--10, 33--35]{edo-3-3-3-jfrp-2026}.

\subsection{Five Core Components}
\begin{enumerate}
  \item \textbf{Reoriented JROC.} Focus on Joint Force Design (the future force), Joint Capability Integration (interdependencies and interoperability), and Combatant Command requirements.
  \item \textbf{Joint Operational Problems.} The \ac{jroc} annually identifies and ranks \acp{jop}, providing the authoritative joint demand signal for analysis, experimentation, and budgeting.
  \item \textbf{Unified urgent needs.} The \ac{jicr} replaces the separate joint urgent and emergent pathways.
  \item \textbf{Empowered Services.} Services generate, validate, document, and approve their requirements while sharing information for joint integration.
  \item \textbf{Direct link to resourcing.} The \ac{rrab} aligns fiscal resources to \ac{jroc}-prioritized problems.
\end{enumerate}
\noindent Source: coursebook slides 9--10~\autocite{edo-3-3-3-jfrp-2026}.

\subsection{Capability Gap Assessment and the Warfighter Voice}
\begin{description}
  \item[\textbf{Annual assessment.}] The \ac{cga} collects, consolidates, assesses, and prioritizes risk-informed \ac{ccmd} capability gaps in an iterative annual process.
  \item[\textbf{Primary input.}] \acp{ccdr} submit \acp{ipl}: prioritized, singular, discrete capability gaps across Service and functional lines that affect assigned missions.
  \item[\textbf{FYDP connection.}] The assessment evaluates whether the current \ac{fydp} addresses those gaps and informs future FYDP recommendations and the Chairman's reporting responsibilities.
  \item[\textbf{CBA crosswalk.}] Older material uses \ac{cba} for a mission/gap/solution study. Know that analytical concept, but answer the new coursebook's annual joint gap-prioritization question with \ac{cga} and \acp{ipl}; the terms are not a universal one-for-one rename.
\end{description}
\noindent Source: coursebook slides 13--15~\autocite{edo-3-3-3-jfrp-2026}.

\subsection{Materiel and Non-Materiel Solutions}
\begin{description}
  \item[\textbf{Non-materiel first.}] Explore doctrine, organization, training, leadership and education, personnel, facilities, and policy changes before assuming a new system is needed. A \ac{dcr} documents a non-materiel approach; these changes can be faster and less costly.
  \item[\textbf{Materiel solution.}] A new or modified item addresses the gap through acquisition. An initial \ac{srd}/\ac{icd} supports further analysis; the \ac{mdauthority}'s \ac{mdd} authorizes entry into the acquisition process.
  \item[\textbf{Combined solutions.}] A capability gap may require both materiel and non-materiel actions. \ac{dotmlpfp} means Doctrine, Organization, Training, Materiel, Leadership and education, Personnel, Facilities, and Policy.
\end{description}
\noindent Source: coursebook slides 14--18~\autocite{edo-3-3-3-jfrp-2026}.
\note{Source ambiguity: slide 14 places a definition of a defense acquisition program under ``Non-materiel Solution.'' Read it with the doctrine/policy examples and the DCR explanation on slides 15--16; do not memorize acquisition of a new system as a non-materiel solution.}

\subsection{Service Requirements Documents and Acquisition Phases}
\begin{longtblr}[
  caption = {Requirements documents in the Major Capability Acquisition teaching flow},
  label = {tab:jfrp-documents},
  entry = {JFRP Requirements Documents},
  remark{\tSource} = {\tabCite{edo-3-3-3-jfrp-2026,edo-3-3-2-acq-milestones-phases-2025}}
]{colspec = {@{} X[1.1] X[1.7] X[1.5] @{}}}
\toprule
{Document} & {Purpose and content} & {Decision / phase linkage} \\
\midrule
Initial SRD / ICD & Describes capability gaps, timeframe, and DOTMLPF-P analysis; no KPPs in the coursebook flow & Approved before MDD; supports MSA, AoA, acquisition strategy, Milestone A, and TMRR \\
Updated SRD / CDD & Defines affordable, supportable, measurable, testable capabilities and performance attributes, including KPPs & Approved before Milestone B; guides EMD and the technical baseline \\
Updated SRD / updated CDD, if required & Refines production attributes, quantities, cost estimates, and later increments & Supports Milestone C and Production and Deployment \\
\bottomrule
\end{longtblr}

\begin{description}
  \item[\textbf{Naming is tailorable.}] \ac{srd} is the coursebook's generic Service-document term, not a universally mandated replacement form. \ac{icd}/\ac{cdd} names may remain in use. A separate \ac{cpd} is a legacy production-document concept~\autocite[slides 7, 20--23]{edo-3-3-3-jfrp-2026}.
  \item[\textbf{Approval distinction.}] Requirements approval establishes the capability need. The \ac{mdauthority} authorizes acquisition phase entry; Milestone~B normally initiates the acquisition program. An approved document does not itself authorize spending or award a contract~\autocite{edo-3-3-3-jfrp-2026,DoDI5000-85}.
  \item[\textbf{Team effort.}] Operators, requirements sponsors, acquisition, engineering, test, logistics, and resource managers must develop executable requirements together. \ac{aoa} during \ac{msa} informs the choice of solution~\autocite{edo-3-3-3-jfrp-2026}.
\end{description}

\subsection{Joint Governance and the Annual Process}
\begin{description}
  \item[\textbf{JROC.}] Senior joint requirements board, chaired by the \ac{vcjcs} with the Service vice chiefs (including the Vice Chief of Space Operations) and Assistant Commandant of the Marine Corps. It prioritizes joint problems and adjudicates joint capability requirements when necessary.
  \item[\textbf{JCB.}] Chaired by Joint Staff J-8; advises the \ac{jroc} and adjudicates requirements, capability assessments, and documents as delegated.
  \item[\textbf{FCBs.}] Seven \acp{fcb}, aligned to Joint Capability Areas, identify/prioritize joint problems, support force design, and continuously review military capabilities. General/flag officer and O-6 integration groups perform initial portfolio-level reviews.
  \item[\textbf{Warfighter participation.}] \acp{ccdr} identify gaps through \acp{ipl} and participation in functional capability boards; retain the operator's voice through analysis and implementation.
\end{description}
\noindent Source: coursebook slides 26--31, 37~\autocite{edo-3-3-3-jfrp-2026}.

The annual teaching sequence is: prioritize \acp{jop} (problems), develop Joint Force Design (joint vision), conduct continuous Joint Capability Integration (integration), use \ac{cpm} to align related investments and requirements (delegation), then review and respond~\autocite[slide 31]{edo-3-3-3-jfrp-2026}.

\subsection{Requirements and Resourcing Theory of Victory}
\begin{description}
  \item[\textbf{JROC demand signal.}] Annually ranks \acp{jop} and identifies associated concept-required capabilities.
  \item[\textbf{RRAB.}] Aligns fiscal resources to prioritized problems during Program Budget Review; selects problems for deeper analysis and funding consideration. By exception it may recommend changing or terminating Component activities that conflict with joint priorities.
  \item[\textbf{MEIA.}] The \ac{meia} conducts mission engineering, experimentation, assessments, and operational integration. It engages industry and provides evidence about solution effectiveness, integration, and interoperability to inform resourcing.
  \item[\textbf{Joint Acceleration Reserve.}] The reproduced manual excerpt directs \ac{cape} to implement the \ac{jar} beginning with the FY~2027 budget cycle: a deliberate topline reserve placing funding into Service execution lines across the \ac{fydp}, with internal tags for visibility. It is not a substitute for appropriations or fiscal controls.
  \item[\textbf{Service execution.}] Services share approved requirements and program data while executing without routine joint-validation delays. Strategy drives priorities, priorities drive resources, and resources drive capability delivery.
\end{description}
\noindent Source: CJCSM 5123.01A, 5 August 2026, excerpts reproduced on coursebook slides 33--35~\autocite{edo-3-3-3-jfrp-2026}.

\subsection{Joint Immediate Combatant Command Requirement}
\begin{description}
  \item[\textbf{Replacement and threshold.}] The \ac{jicr} consolidates \ac{juon}/\ac{jeon}. It addresses an immediate joint \ac{ccmd} gap that, unfulfilled, will result in certain death and/or unacceptable risk to a primary objective of ongoing or imminent operations~\autocite[slide 10]{edo-3-3-3-jfrp-2026}.
  \item[\textbf{Timing.}] The coursebook describes a 35-business-day assessment/approval process and a goal of fielding usable capability within two years. This is not a 35-day fielding promise~\autocite[slide 10]{edo-3-3-3-jfrp-2026}.
  \item[\textbf{Execution.}] The Urgent Capability Acquisition pathway is the acquisition mechanism for validated urgent needs; requirements validation and acquisition execution are separate decisions~\autocite{edo-3-3-2-acq-milestones-phases-2025}.
  \item[\textbf{Disposition.}] Within six months of fielding, the requirements sponsor submits an Assessment of Operational Utility to the solution sponsor and approval authority. It informs whether to transition to a program of record, sustain for the current contingency, or cancel~\autocite{edo-3-3-2-acq-milestones-phases-2025}.
\end{description}

\subsection{Performance Attributes and Test Evidence}
\begin{longtblr}[
  caption = {KPP, KSA, and APA distinctions},
  label = {tab:kpp-ksa-apa-comparison},
  entry = {KPP vs KSA vs APA Comparison},
  remark{\tSource} = {\tabCite{DoDI5000-85,edo-3-3-3-jfrp-2026}}
]{colspec = {@{} X[0.65] X[1.7] X[2] @{}}}
\toprule
{Attribute} & {Meaning} & {Decision consequence} \\
\midrule
KPP & Critical or essential performance attribute & A threshold miss can threaten operational utility and APB compliance; assess corrective action, baseline implications, and decision-authority action \\
KSA & Important to a balanced solution or operational suitability; below KPP criticality & A miss still creates capability, suitability, integration, or acceptance risk \\
APA & Additional attribute needed for full capability; below KPP/KSA criticality & A miss still requires test evidence, impact assessment, corrective action, and explicit residual-risk disposition \\
\bottomrule
\end{longtblr}
\noindent All attributes must be measurable and testable. Trace approved requirements through engineering allocations to verification methods, the \ac{temp}, and milestone evidence. \acp{kpp} are carried into the \ac{apb}; do not assume every KSA or APA is automatically an APB performance parameter, or that a KPP miss automatically triggers a particular statutory report~\autocite{DoDI5000-85,edo-3-3-3-jfrp-2026}.

\subsection{Information Technology and Interoperability}
\begin{description}
  \item[\textbf{Net-Ready.}] The coursebook documents interoperability needs as a Net-Ready \ac{kpp} or \ac{ksa}: IT supports military operations, can enter and be managed on the network, and effectively exchanges information.
  \item[\textbf{Architecture and reuse.}] \ac{dodaf} defines architecture views and encourages reuse; technical views must be consistent with \ac{disr}. The coursebook identifies architecture views as required parts of a \ac{cdd}.
  \item[\textbf{Joint and allied operations.}] Consider all affected U.S. and allied systems. Stovepipes waste resources; failed interoperability can compromise operations and put lives and equipment at risk.
\end{description}
\noindent Source: coursebook slides 28--29~\autocite{edo-3-3-3-jfrp-2026}.

\subsection{Quick Review}
\begin{itemize}
  \item What replaced JCIDS? \textbf{Answer:} \ac{jfrp}; Services approve their requirements while the \ac{jroc} focuses on joint problems and integration~\autocite{edo-3-3-3-jfrp-2026}.
  \item Name the five core components. \textbf{Answer:} Reoriented JROC, prioritized JOPs, unified JICR process, empowered Services, and direct requirements-to-resourcing alignment~\autocite{edo-3-3-3-jfrp-2026}.
  \item What feeds the annual CGA? \textbf{Answer:} Combatant Commander Integrated Priority Lists~\autocite{edo-3-3-3-jfrp-2026}.
  \item Which documents support MDD and Milestone B? \textbf{Answer:} Initial SRD/ICD before MDD; updated SRD/CDD before Milestone~B. The MDA approves acquisition phase entry~\autocite{edo-3-3-3-jfrp-2026}.
  \item What replaced JUON/JEON, and does approval mean immediate fielding? \textbf{Answer:} JICR; the coursebook distinguishes a 35-business-day assessment/approval process from fielding within two years~\autocite{edo-3-3-3-jfrp-2026}.
  \item What do RRAB and MEIA do? \textbf{Answer:} RRAB aligns resources to joint priorities; MEIA provides mission-engineering, experimentation, and integration evidence~\autocite[slides 33--35]{edo-3-3-3-jfrp-2026}.
\end{itemize}

%====================
% End of file
%====================
\IfSubfilesClassLoaded{
  \RenewDocumentCommand{\entryname}{}{\textbf{\color{Modern} Acronym}}
  \RenewDocumentCommand{\descriptionname}{}{\textbf{\color{Modern} Definition}}
  \printnoidxglossary[
    type=\acronymtype,
    title=Acronyms,
    style=long-booktabs
  ]}{}

\end{document}
